\newpage\thispagestyle{hl}
\hlbanner{Part III $\cdot$ Reduced-order modelling}{12}{Non-intrusive identification and nonlinear ROMs}
\hlsection{The question}
\begin{hlquote}A vendor gives you a black-box solver for a structure that bends a lot. You can run it, but you cannot see its matrices. Everything so far projected matrices; here there are none to project.\end{hlquote}
\vspace{8pt}\begin{hlsum}\textbf{When only a black-box solver is available, the reduced model is identified from its outputs.} At a mid-span deflection of \cnlDefl{} mm the nonlinear effect is \cnlNlEffectPct\%. With 3 modes a linear model errs by \cnlLinThree, a polynomial model by \cnlIceThree{} and an RBF surrogate by \cnlMfsThree. The chapter identifies modal frequencies from a frequency response, tests nonlinear force models against full-order solves not used in training, and traces a backbone by harmonic balance. Loewner identification proved fragile: a fifth of the random point choices missed a mode even without noise.\end{hlsum}
\hlsection{Key concepts}
\begin{description}\setlength\itemsep{5pt}\setlength\parsep{0pt}\renewcommand{\descriptionlabel}[1]{\hspace{\labelsep}\sffamily\bfseries\color{hlnavy}#1}
\item[Non-intrusive] models built from outputs only, without access to matrices
\item[Loewner identification] natural frequencies and damping from sampled frequency responses
\item[Operator inference] reduced mass and stiffness learned by least squares, kept positive definite
\item[Hybrid nonlinear modal model] linear part intrusive, nonlinear force fitted by polynomial or RBF
\item[Nonlinear normal mode] backbone of frequency against amplitude, by harmonic balance
\item[Hardening] large deflection stiffens a clamped beam through a cubic term
\end{description}
\hlsection{After this chapter you can}
\begin{itemize}\setlength\itemsep{3pt}
\item identify modal frequencies from a frequency response and know when to distrust the result
\item build and test a polynomial or RBF modal force model against full-order solves not used in training
\item trace and check a backbone by harmonic balance
\end{itemize}
\hlsection{Explained on the page, in short boxes}
\noindent\small Intrusive and non-intrusive; Identification from frequency response data; Why large deflection stiffens a beam; Backbone.
\newpage\thispagestyle{hl}
\noindent{\sffamily\small\color{hlgrey}Chapter 12 \quad Non-intrusive identification and nonlinear ROMs}\par
\hlfig{\hlimg{ch12_nlrom}}{Mid-span deflection of a clamped--clamped beam against load for the full-order and reduced models (left), and largest error against retained modes (right).}
\hlsection{Headline result}
\begin{hlheadline}At a mid-span deflection of \cnlDefl{} mm the nonlinear effect is \cnlNlEffectPct\%. With 3 modes a linear model errs by \cnlLinThree, the polynomial model by \cnlIceThree{} and the RBF surrogate by \cnlMfsThree.\end{hlheadline}
\hlsection{Main classes and functions}
{\begin{itemize}\setlength\itemsep{2pt}
\item \texttt{rom\_engine.nonlinear\_rom.PolynomialModalROM}
\item \texttt{rom\_engine.nonlinear\_rom.MultiFidelitySurrogate}
\item \texttt{rom\_engine.nnm.solve\_nnm\_backbone}
\item \texttt{rom\_engine.operator\_inference.infer\_contact\_operators}
\end{itemize}}
\hlsection{A stated limit}
\begin{hllimit}Loewner identification is fragile: even with noise-free data a fifth of the random choices of interpolation points missed a mode, and no draw succeeded at 10\% noise.\end{hllimit}
